\documentclass[../main]{subfiles}
\begin{document}
\chapter{Pérdidas en Motores.}
\begin{center}
  \img{images/09/placamotor.jpg}{Placa de características de motor}
\end{center}

\info{¿Qué son las pérdidas?}{Las pérdidas de los motores eléctricos se producen en las láminas del estator, las láminas del rotor y los devanados, pero también en los imanes permanentes. Cuanto menores sean las pérdidas del motor eléctrico, mayor será su eficiencia.}

\section{Cálculo de pérdidas}

\subsection{Pérdidas en el estator}

\subsubsection{Potencia eléctrica en la entrada del estator}
\begin{equation}
  P_{in} = m \cdot V \cdot I \cdot \cos(\varphi)
\end{equation}

Siendo:
\begin{itemize}
  \item $P_{in}$: Potencia de entrada en el estator (potencia eléctrica absorbida de la red) \unit{W}.
  \item $m$: número de fases del estator.
  \item $V$: Tensión por fase del estator \unit{V}.
  \item $I$: Corriente del estator por fase \unit{A}.
  \item $\cos(\varphi)$: factor de potencia del motor (f.d.p del estator).
\end{itemize}

\subsubsection{Pérdidas en el cobre del estator ($P_{Cu1}$)}
\begin{equation}
  P_{Cu1} = m \cdot R_1 \cdot I_1^2
\end{equation}

\subsubsection{Pérdidas totales en el estator}
\begin{equation}
  P_{p} = P_{Cu1} + P_{Fe}
\end{equation}

\subsection{Pérdidas en el hierro del estator}
\begin{equation}
  P_{Fe} = m \cdot E_1 \cdot I_{Fe}
\end{equation}
Donde $I_{Fe}$ es la Corriente de pérdidas en el hierro.

\subsection{Potencia que atraviesa el entrehierro}
\begin{equation}
  P_a = P_{in} - P_p = P_{in} - P_{Cu1} - P_{Fe}
\end{equation}

\subsection{Pérdidas en el rotor}

\subsubsection{Pérdidas en el cobre del rotor}
\begin{equation}
  P_{Cu2} = m \cdot R_2 \cdot I_2^2
\end{equation}

\subsection{Potencia mecánica interna del motor}
\begin{equation}
  P_m = P_a - P_{Cu2}
\end{equation}

\subsection{Expresión de la potencia mecánica interna}
\begin{equation}
  P_m = m \cdot R_2 (1-s) I_2^2
\end{equation}

\subsection{Potencia mecánica útil del motor}
\begin{equation}
  P_u = P_i - P_m
\end{equation}

\subsection{Rendimiento del motor}
\begin{equation}
  \eta = \dfrac{P_u}{P_{in}} = \dfrac{P_u}{P_u + P_m + P_{Cu2} + P_{Fe} + P_{Cu1}}
\end{equation}

\actividad{Problemas}
\begin{enumerate}
  \item La potencia absorbida por un motor asíncrono trifásico de 4 polos, \qty{50}{Hz}, es de \qty{476}{kW} cuando gira a \qty{1435}{rpm}. Las pérdidas totales en el estator son de \qty{265}{W} y las de rozamiento y ventilación son de \qty{300}{W}. Calcular:
    \begin{enumerate}
      \item el deslizamiento,
      \item las pérdidas en el rotor,
      \item potencia útil en el árbol del motor,
      \item rendimiento.
    \end{enumerate}
  \item Un motor de inducción trifásico de 8 polos, \qty{1}{CV}, \qty{380}{V} gira a \qty{72}{rpm} a plena carga. Si el rendimiento y el f.d.p. a esta carga es de 83\% y 0,75 respectivamente. Calcular:
    \begin{enumerate}
      \item velocidad de sincronismo del campo giratorio,
      \item deslizamiento a plena carga,
      \item corriente de línea.
    \end{enumerate}
  \item Un motor asincrónico trifásico de 4 polos, conectado en estrella se alimenta por una red de \qty{380}{V}, \qty{50}{Hz}. La impedancia del estator es igual a $Z_1 = 0,1 + j0,5~\Omega$ y la del rotor en reposo reducida al estator vale $Z_2 = 0,1 + j0,3~\Omega/\text{fase}$. Calcular:
    \begin{enumerate}
      \item intensidad absorbida en el arranque,
      \item corriente a plena carga, si el deslizamiento es de 4\%,
      \item potencia y par nominal si se desprecian las pérdidas mecánicas,
      \item rendimiento en el caso anterior si las pérdidas en el hierro son iguales a \qty{1200}{W}.
    \end{enumerate}
  \item Un motor de inducción trifásico de 6 polos, \qty{60}{Hz}, absorbe una potencia de \qty{200}{kW} cuando gira a \qty{1180}{rpm}. Las pérdidas en el estator son de \qty{180}{W} y las pérdidas por fricción son de \qty{250}{W}. Calcular:
    \begin{enumerate}
      \item el deslizamiento,
      \item las pérdidas en el rotor,
      \item la potencia útil en el eje del motor,
      \item el rendimiento.
    \end{enumerate}
  \item Un motor asincrónico de 10 polos, \qty{1.5}{CV}, \qty{440}{V}, opera a \qty{58}{rpm} a plena carga. El rendimiento es del 85\% y el factor de potencia es de 0.78. Calcular:
    \begin{enumerate}
      \item la velocidad sincrónica,
      \item el deslizamiento a plena carga,
      \item la corriente de línea,
      \item la potencia entregada en el eje del motor.
    \end{enumerate}
  \item Un motor de inducción trifásico de 4 polos, \qty{50}{Hz}, conectado en triángulo a una red de \qty{400}{V} tiene una impedancia del estator de $Z_1 = 0.12 + j0.6~\Omega$. La impedancia del rotor en reposo es $Z_2 = 0.15 + j0.35~\Omega$. Calcular:
    \begin{enumerate}
      \item la corriente absorbida en el arranque,
      \item la corriente a plena carga con un deslizamiento del 3\%,
      \item el par nominal suponiendo despreciables las pérdidas mecánicas,
      \item el rendimiento si las pérdidas en el hierro son de \qty{1500}{W}.
    \end{enumerate}
  \item Un motor asíncrono de 8 polos, \qty{1}{CV}, \qty{380}{V}, tiene un deslizamiento de 2\% cuando gira a plena carga. Las pérdidas totales en el estator y rotor son de \qty{500}{W}. Calcular:
    \begin{enumerate}
      \item la velocidad sincrónica,
      \item la velocidad del rotor a plena carga,
      \item la potencia en el eje del motor,
      \item el rendimiento.
    \end{enumerate}
  \item Un motor de inducción de 6 polos, \qty{60}{Hz}, con una potencia absorbida de \qty{100}{kW}, gira a \qty{1150}{rpm}. Las pérdidas en el estator son de \qty{300}{W} y las pérdidas mecánicas son de \qty{350}{W}. Calcular:
    \begin{enumerate}
      \item el deslizamiento,
      \item las pérdidas en el rotor,
      \item la potencia útil en el eje del motor,
      \item el rendimiento.
    \end{enumerate}
  \item Un motor asincrónico trifásico de 10 polos, \qty{2}{CV}, \qty{400}{V}, conectado en estrella, gira a \qty{68}{rpm}. El rendimiento a plena carga es del 80\% y el factor de potencia es de 0.8. Calcular:
    \begin{enumerate}
      \item la velocidad sincrónica,
      \item el deslizamiento a plena carga,
      \item la corriente de línea,
      \item la potencia entregada en el eje del motor.
    \end{enumerate}
  \item Un motor de inducción de 4 polos, \qty{380}{V}, \qty{50}{Hz} tiene una impedancia del estator de $Z_1 = 0.1 + j0.4~\Omega$. La impedancia del rotor en reposo es de $Z_2 = 0.12 + j0.3~\Omega$. El deslizamiento a plena carga es del 5\%. Calcular:
    \begin{enumerate}
      \item la corriente en el arranque,
      \item la corriente a plena carga,
      \item el par nominal si las pérdidas mecánicas son despreciables,
      \item el rendimiento si las pérdidas en el hierro son de \qty{1000}{W}.
    \end{enumerate}
  \item Un motor asincrónico trifásico de 6 polos, conectado en estrella, con una potencia de \qty{150}{kW}, gira a \qty{1200}{rpm}. Las pérdidas en el rotor son de \qty{500}{W} y las pérdidas en el estator son de \qty{400}{W}. Calcular:
    \begin{enumerate}
      \item el deslizamiento,
      \item la potencia útil en el eje del motor,
      \item el rendimiento,
      \item el par desarrollado.
    \end{enumerate}
  \item Un motor de inducción trifásico de 4 polos, \qty{440}{V}, \qty{60}{Hz}, conectado en triángulo, tiene una impedancia del estator de $Z_1 = 0.08 + j0.5~\Omega$. La impedancia del rotor en reposo es de $Z_2 = 0.09 + j0.35~\Omega$. Calcular:
    \begin{enumerate}
      \item la corriente absorbida en el arranque,
      \item la corriente a plena carga con un deslizamiento del 2\%,
      \item la potencia nominal si las pérdidas mecánicas son despreciables,
      \item el rendimiento si las pérdidas en el hierro son de \qty{1200}{W}.
    \end{enumerate}
  \item Un motor asincrónico trifásico de 8 polos, \qty{2}{CV}, \qty{380}{V}, tiene un deslizamiento de 1.5\% a plena carga. Las pérdidas en el estator y el rotor suman \qty{600}{W}. Calcular:
    \begin{enumerate}
      \item la velocidad sincrónica,
      \item la velocidad del rotor a plena carga,
      \item la potencia entregada en el eje del motor,
      \item el rendimiento.
    \end{enumerate}
\end{enumerate}

\end{document}